\documentclass[11pt]{article}
\usepackage[margin=1in]{geometry}
\usepackage[T1]{fontenc}
\usepackage{lmodern}
\usepackage{amsmath,amssymb}
\usepackage{hyperref}

\title{Literature Status for arXiv:math/0406477: Exact Complexity}
\author{Run fa\_banach\_001, agent\_lane\_12}
\date{2026-06-16}

\begin{document}
\maketitle

\section*{Status}

This is a \texttt{literature\_already\_answered} packet.  It records that the
exact-complexity question raised in Remark~5.4 of Ferenczi--Medina Galego,
\emph{Some equivalence relations which are Borel reducible to isomorphism
between separable Banach spaces}, arXiv:math/0406477, was answered by
Ferenczi--Louveau--Rosendal, \emph{The complexity of classifying separable
Banach spaces up to isomorphism}, arXiv:math/0610289.

\section*{Source Question}

Remark~5.4 of arXiv:math/0406477 says that the exact complexity of isomorphism
between separable Banach spaces, viewed as subspaces of \(C[0,1]\) with the
Effros--Borel structure, was ``quite open''.  It also says the same was true
``even'' for Banach spaces with Schauder bases, viewed as subsequences of the
universal basis of Pelczynski, and suggests the possibility that the relation
could be \(\Sigma^1_1\)-complete, i.e. maximum among analytic equivalence
relations under Borel reducibility.

\section*{Supporting Answer}

The abstract of arXiv:math/0610289 states that isomorphism between separable
Banach spaces is a complete analytic equivalence relation, meaning that every
analytic equivalence relation Borel reduces to it.  The main theorem in the
section on isomorphism of separable Banach spaces repeats this statement:
\[
\text{isomorphism between separable Banach spaces is complete analytic.}
\]
This is exactly the \(\Sigma^1_1\)-complete alternative raised in the source
remark.

The supporting paper also addresses the basis-side concern from the same
remark.  Theorem~1 proves that permutative equivalence between, even
unconditional, Schauder bases is a complete analytic equivalence relation.
Moreover, the later Banach-space reduction uses the spaces
\(\Delta(Z_S,2)\), which are reflexive and carry suppression unconditional
bases.  Thus the later theorem answers the general separable Banach-space
exact-complexity question and supplies the expected basis-level complete
analytic phenomenon.

\section*{Scope}

This packet does not address the separate questions in Remark~5.2 and
Conjecture~5.3 of arXiv:math/0406477: ergodicity of \(\ell_p\) for \(p>2\),
the conjecture that \(\ell_2\) is the only non-ergodic Banach space, and the
block-subspace dichotomy for Banach spaces with unconditional bases remain
outside this literature-status claim.

\section*{Search Evidence}

The lane-12 queue selected arXiv:math/0406477.  Cheap run indexes were searched
for the arXiv id, title words, Ferenczi--Galego, Borel reducibility, and
exact-complexity phrases; no duplicate packet for this source was found.
Targeted web and local searches for phrases including
\texttt{isomorphism between separable Banach spaces complete analytic},
\texttt{Sigma\^{}1\_1-complete}, \texttt{Schauder bases}, and
\texttt{universal basis of Pelczynski} identified arXiv:math/0610289 as the
decisive later paper.

\begin{thebibliography}{9}
\bibitem{FG04}
V.~Ferenczi and E.~Medina Galego,
\emph{Some equivalence relations which are Borel reducible to isomorphism
between separable Banach spaces},
arXiv:math/0406477.

\bibitem{FLR06}
V.~Ferenczi, A.~Louveau, and C.~Rosendal,
\emph{The complexity of classifying separable Banach spaces up to isomorphism},
arXiv:math/0610289.
\end{thebibliography}

\end{document}
